\nonstopmode{}
\documentclass[a4paper]{book}
\usepackage[times,inconsolata,hyper]{Rd}
\usepackage{makeidx}
\makeatletter\@ifl@t@r\fmtversion{2018/04/01}{}{\usepackage[utf8]{inputenc}}\makeatother
% \usepackage{graphicx} % @USE GRAPHICX@
\makeindex{}
\begin{document}
\chapter*{}
\begin{center}
{\textbf{\huge Package `crypto'}}
\par\bigskip{\large \today}
\end{center}
\ifthenelse{\boolean{Rd@use@hyper}}{\hypersetup{pdftitle = {crypto: Interactive Cryptocurrency Time Series and Market Data Explorer}}}{}
\ifthenelse{\boolean{Rd@use@hyper}}{\hypersetup{pdfauthor = {Dimitrios Zacharatos}}}{}
\begin{description}
\raggedright{}
\item[Title]\AsIs{Interactive Cryptocurrency Time Series and Market Data Explorer}
\item[Version]\AsIs{0.1.0}
\item[Maintainer]\AsIs{Dimitrios Zacharatos }\email{dzacharatos@yahoo.com}\AsIs{}
\item[Description]\AsIs{A 'shiny' application for exploring cryptocurrency exchange rates.
Plots interactive candlestick and volume charts of historical prices since
2013 for any active cryptocurrency, and shows a table of the top 100
cryptocurrencies by market capitalisation with price changes over 24 hours
to 90 days, market dominance and fully diluted market cap. Also compares
normalized prices of several coins, and shows a market cap treemap and a
calendar of monthly returns. Market data is retrieved from
'CoinMarketCap'.}
\item[Imports]\AsIs{bslib,
cachem,
coinmarketcapr,
crypto2,
DT,
plotly,
shiny,
shinycssloaders,
shinyWidgets,
stats,
tools,
utils}
\item[Suggests]\AsIs{testthat (>= 3.0.0)}
\item[Config/testthat/edition]\AsIs{3}
\item[License]\AsIs{GPL-3}
\item[URL]\AsIs{}\url{https://sedzinfo.github.io/crypto/}\AsIs{, }\url{https://github.com/sedzinfo/crypto}\AsIs{}
\item[BugReports]\AsIs{}\url{https://github.com/sedzinfo/crypto/issues}\AsIs{}
\item[Encoding]\AsIs{UTF-8}
\item[Roxygen]\AsIs{list(markdown = TRUE)}
\item[Config/roxygen2/version]\AsIs{8.1.1}
\end{description}
\Rdcontents{Contents}
\HeaderA{crypto}{Launch the Crypto Time Series Analyzer}{crypto}
\keyword{crypto}{crypto}
\keyword{timeseries}{crypto}
%
\begin{Description}
Starts an interactive 'shiny' application for exploring
cryptocurrency exchange rates and market data from 'CoinMarketCap'.
\end{Description}
%
\begin{Usage}
\begin{verbatim}
crypto(api_key = Sys.getenv("COINMARKETCAP_APIKEY"), ...)
\end{verbatim}
\end{Usage}
%
\begin{Arguments}
\begin{ldescription}
\item[\code{api\_key}] A 'CoinMarketCap' API key, needed for the market listings.
By default it is read from the \code{COINMARKETCAP\_APIKEY} environment
variable, which can be set in \code{.Renviron}.

\item[\code{...}] Further arguments passed to \code{\LinkA{shiny::shinyApp()}{shiny::shinyApp()}}, for example
\code{options = list(port = 8080)}.
\end{ldescription}
\end{Arguments}
%
\begin{Details}
The application has the following tabs:
\begin{description}

\item[Timeseries] Candlestick chart of daily open, high, low and close
prices in USD since 2013 for any active cryptocurrency, with its 50-
and 200-day moving averages, and charts of the volume and of the
drawdown (how far the price is below its all-time high) underneath.
\item[Top 100] Table of the top 100 cryptocurrencies by market cap with
price, volume, 24h to 90d price changes, market dominance and fully
diluted market cap.
\item[Compare] Close prices of up to 10 coins, normalized to 100 on the
first date all of them have data.
\item[Treemap] Top 100 coins sized by market cap and colored by 24h
change.
\item[Returns calendar] Monthly returns by year for the coin selected
in the Timeseries tab.

\end{description}

Market listings are downloaded when the app starts and refreshed every 10
minutes. Historical prices are downloaded on demand with
\code{\LinkA{crypto2::crypto\_history()}{crypto2::crypto.Rul.history()}} and kept on disk between R sessions, so
after the first download only the days since the last one are fetched.
An internet connection is required. When 'CoinMarketCap' limits the requests,
crypto2 prints "Retrying in 60 seconds" and the chart appears about a
minute later.
\end{Details}
%
\begin{Value}
A \code{\LinkA{shiny::shinyApp()}{shiny::shinyApp()}} object. Printing it, or calling \code{crypto()} at
the console, opens the app in the browser.
\end{Value}
%
\begin{Examples}
\begin{ExampleCode}
if (interactive()) crypto()

# with the API key passed directly, on a fixed port
if (interactive()) crypto(api_key = "your-api-key", options = list(port = 8080))
\end{ExampleCode}
\end{Examples}
\HeaderA{drawdown}{Drawdown of prices}{drawdown}
%
\begin{Description}
How far each close is below the highest close up to that day, the
all-time high so far, in percent: 0 on a new high and -50 at half of it.
\end{Description}
%
\begin{Usage}
\begin{verbatim}
drawdown(close)
\end{verbatim}
\end{Usage}
%
\begin{Arguments}
\begin{ldescription}
\item[\code{close}] Numeric vector of close prices, sorted by date.
\end{ldescription}
\end{Arguments}
%
\begin{Value}
A numeric vector as long as \code{close}, from -100 to 0.
\end{Value}
%
\begin{Examples}
\begin{ExampleCode}
drawdown(c(100, 120, 90, 130, 65))
\end{ExampleCode}
\end{Examples}
\HeaderA{get\_coins}{List the active coins}{get.Rul.coins}
%
\begin{Description}
Downloads the list of the active coins tracked by 'CoinMarketCap' and
keeps it on disk for a day (see \code{\LinkA{disk\_cache()}{disk.Rul.cache}}). When the download fails,
the kept list is used. Its \code{name\_symbol} column holds the labels, such as
"Bitcoin (BTC)", that pick coins in \code{\LinkA{get\_history()}{get.Rul.history}}.
\end{Description}
%
\begin{Usage}
\begin{verbatim}
get_coins()
\end{verbatim}
\end{Usage}
%
\begin{Value}
The \code{\LinkA{crypto2::crypto\_list()}{crypto2::crypto.Rul.list()}} data frame, with a \code{name\_symbol}
column.
\end{Value}
%
\begin{Examples}
\begin{ExampleCode}
## Not run: 
coins <- get_coins()
head(coins$name_symbol)

## End(Not run)
\end{ExampleCode}
\end{Examples}
\HeaderA{get\_history}{Get the daily price history of coins}{get.Rul.history}
%
\begin{Description}
Returns the daily USD prices of one or more coins from 'CoinMarketCap'.
Each coin's history is kept on disk between R sessions, and once it is
more than an hour old only the days since its last day are downloaded, so
calls are fast and stay below the request limit of 'CoinMarketCap'. When
that limit is hit, crypto2 waits a minute and tries again.
\end{Description}
%
\begin{Usage}
\begin{verbatim}
get_history(labels, start_date = Sys.Date() - 365, coins = get_coins())
\end{verbatim}
\end{Usage}
%
\begin{Arguments}
\begin{ldescription}
\item[\code{labels}] Character vector of coin labels, such as "Bitcoin (BTC)",
from the \code{name\_symbol} column of \code{\LinkA{get\_coins()}{get.Rul.coins}}.

\item[\code{start\_date}] Date. The first day to get, a year ago by default.

\item[\code{coins}] Data frame of coins from \code{\LinkA{get\_coins()}{get.Rul.coins}}, to look the labels up.
\end{ldescription}
\end{Arguments}
%
\begin{Value}
A data frame with one row per coin and day, sorted by coin and
date, with the \code{open}, \code{high}, \code{low}, \code{close} and \code{volume} columns,
\code{timestamp} as a Date and the \code{name\_symbol} label. Labels that match no
coin are left out with a message, and \code{NULL} is returned when no label
matches or nothing could be downloaded.
\end{Value}
%
\begin{Examples}
\begin{ExampleCode}
## Not run: 
btc <- get_history("Bitcoin (BTC)", as.Date("2020-01-01"))
plot_timeseries(btc)

history <- get_history(c("Bitcoin (BTC)", "Ethereum (ETH)"))
plot_compare(normalize_prices(history))

## End(Not run)
\end{ExampleCode}
\end{Examples}
\HeaderA{get\_listings}{Get the latest market listings}{get.Rul.listings}
%
\begin{Description}
Returns the listings of the coins by market cap, with their price, volume,
price changes from 24 hours to 90 days, market cap and dominance in USD.
The listings are downloaded again when they are older than 10 minutes,
and the download is shared by all sessions and app runs in the same R
process, so the API is called at most once in that time.
\end{Description}
%
\begin{Usage}
\begin{verbatim}
get_listings()
\end{verbatim}
\end{Usage}
%
\begin{Details}
The listings need a 'CoinMarketCap' API key in the \code{COINMARKETCAP\_APIKEY}
environment variable, see the README. The first download lets errors
through. When a later download fails, the previous listings and their
time are kept, so the time stays truthful.
\end{Details}
%
\begin{Value}
A list with \code{data}, a data frame with one row per coin sorted by
market cap, and \code{time}, when it was downloaded.
\end{Value}
%
\begin{Examples}
\begin{ExampleCode}
## Not run: 
listings <- get_listings()
listings$time
head(listings$data)

## End(Not run)
\end{ExampleCode}
\end{Examples}
\HeaderA{listings\_table}{Table of the market listings}{listings.Rul.table}
%
\begin{Description}
A sortable table with 100 rows per page and numbers rounded to 2
decimals, showing the price in dollars and the price changes in percent.
Losses are shaded red and gains green, more strongly the larger the
change, in steps of 0.1 percentile from the 5th to the 95th. Changes near
zero keep the table's own background, so the table suits the light and the
dark mode.
\end{Description}
%
\begin{Usage}
\begin{verbatim}
listings_table(listings)
\end{verbatim}
\end{Usage}
%
\begin{Arguments}
\begin{ldescription}
\item[\code{listings}] The \code{data} of \code{\LinkA{get\_listings()}{get.Rul.listings}}.
\end{ldescription}
\end{Arguments}
%
\begin{Value}
A DT datatable.
\end{Value}
%
\begin{Examples}
\begin{ExampleCode}
## Not run: 
listings_table(get_listings()$data)

## End(Not run)
\end{ExampleCode}
\end{Examples}
\HeaderA{monthly\_returns}{Monthly returns of a coin}{monthly.Rul.returns}
%
\begin{Description}
Each month's return is its last close relative to the last close of the
month before, in percent. The first month, and the months still to come,
are \code{NA}.
\end{Description}
%
\begin{Usage}
\begin{verbatim}
monthly_returns(history)
\end{verbatim}
\end{Usage}
%
\begin{Arguments}
\begin{ldescription}
\item[\code{history}] Price history of one coin from \code{\LinkA{get\_history()}{get.Rul.history}}, with
\code{timestamp} and \code{close} columns, sorted by date.
\end{ldescription}
\end{Arguments}
%
\begin{Value}
A year by month matrix of returns in percent, rounded to 2
decimals, with the years as row names and \code{month.abb} as column names.
\end{Value}
%
\begin{Examples}
\begin{ExampleCode}
history <- data.frame(
  timestamp = as.Date(c("2024-01-31", "2024-02-29", "2024-03-31")),
  close = c(100, 110, 121)
)
monthly_returns(history)
\end{ExampleCode}
\end{Examples}
\HeaderA{moving\_average}{Moving average of prices}{moving.Rul.average}
%
\begin{Description}
The simple moving average: each day's value is the mean of the last
\code{days} closes, that day included, so the first \code{days - 1} days have none.
\end{Description}
%
\begin{Usage}
\begin{verbatim}
moving_average(close, days)
\end{verbatim}
\end{Usage}
%
\begin{Arguments}
\begin{ldescription}
\item[\code{close}] Numeric vector of close prices, sorted by date.

\item[\code{days}] Number of days to average, such as 50 or 200.
\end{ldescription}
\end{Arguments}
%
\begin{Value}
A numeric vector as long as \code{close}, \code{NA} for the first
\code{days - 1} days, and all \code{NA} when there are fewer than \code{days} closes,
as for a coin younger than the average.
\end{Value}
%
\begin{Examples}
\begin{ExampleCode}
moving_average(c(1, 2, 3, 4, 5), days = 3)
\end{ExampleCode}
\end{Examples}
\HeaderA{normalize\_prices}{Normalize close prices to 100}{normalize.Rul.prices}
%
\begin{Description}
Every coin starts at 100 on the first date on which all of them have data,
so their lines can be compared directly: 150 means the price rose by half
since that date.
\end{Description}
%
\begin{Usage}
\begin{verbatim}
normalize_prices(history)
\end{verbatim}
\end{Usage}
%
\begin{Arguments}
\begin{ldescription}
\item[\code{history}] Price history from \code{\LinkA{get\_history()}{get.Rul.history}}, with \code{name\_symbol},
\code{timestamp} and \code{close} columns, sorted by coin and date.
\end{ldescription}
\end{Arguments}
%
\begin{Value}
\code{history} from the first common date on, with an \code{index} column.
\end{Value}
%
\begin{Examples}
\begin{ExampleCode}
history <- data.frame(
  name_symbol = rep(c("Alpha (A)", "Beta (B)"), each = 3),
  timestamp = rep(as.Date("2024-01-01") + 0:2, 2),
  close = c(10, 11, 12, 200, 180, 220)
)
normalize_prices(history)
\end{ExampleCode}
\end{Examples}
\HeaderA{plot\_compare}{Line chart of normalized close prices}{plot.Rul.compare}
%
\begin{Description}
One line per coin, with buttons that switch the y axis between a linear
and a log scale.
\end{Description}
%
\begin{Usage}
\begin{verbatim}
plot_compare(prices, mode = "light")
\end{verbatim}
\end{Usage}
%
\begin{Arguments}
\begin{ldescription}
\item[\code{prices}] Prices from \code{\LinkA{normalize\_prices()}{normalize.Rul.prices}}.

\item[\code{mode}] \code{"light"} (the default) or \code{"dark"}, to suit a light or dark
page.
\end{ldescription}
\end{Arguments}
%
\begin{Value}
A plotly object.
\end{Value}
%
\begin{Examples}
\begin{ExampleCode}
## Not run: 
history <- get_history(c("Bitcoin (BTC)", "Ethereum (ETH)", "Solana (SOL)"))
plot_compare(normalize_prices(history))

## End(Not run)
\end{ExampleCode}
\end{Examples}
\HeaderA{plot\_returns\_calendar}{Heatmap of monthly returns by year}{plot.Rul.returns.Rul.calendar}
%
\begin{Description}
Months are colored from red to green by their return. The color limit is
the 95th percentile of the absolute returns, so a few extreme months do
not wash out the rest.
\end{Description}
%
\begin{Usage}
\begin{verbatim}
plot_returns_calendar(calendar, label = NULL, mode = "light")
\end{verbatim}
\end{Usage}
%
\begin{Arguments}
\begin{ldescription}
\item[\code{calendar}] Matrix from \code{\LinkA{monthly\_returns()}{monthly.Rul.returns}}.

\item[\code{label}] Label of the coin for the title, or \code{NULL} for none.

\item[\code{mode}] \code{"light"} (the default) or \code{"dark"}, to suit a light or dark
page.
\end{ldescription}
\end{Arguments}
%
\begin{Value}
A plotly object.
\end{Value}
%
\begin{Examples}
\begin{ExampleCode}
## Not run: 
btc <- get_history("Bitcoin (BTC)", as.Date("2013-04-29"))
plot_returns_calendar(monthly_returns(btc), "Bitcoin (BTC)")

## End(Not run)
\end{ExampleCode}
\end{Examples}
\HeaderA{plot\_timeseries}{Candlestick chart of a coin's prices}{plot.Rul.timeseries}
%
\begin{Description}
Plots the daily open, high, low and close prices, with a line through the
close, its moving averages, date range buttons and a range slider.
Underneath, smaller charts show the trading volume and the drawdown: how
far the price is below its all-time high so far. A click on a legend entry
hides or shows that line.
\end{Description}
%
\begin{Usage}
\begin{verbatim}
plot_timeseries(
  history,
  mode = "light",
  averages = c(50, 200),
  show_drawdown = TRUE
)
\end{verbatim}
\end{Usage}
%
\begin{Arguments}
\begin{ldescription}
\item[\code{history}] Price history of one coin from \code{\LinkA{get\_history()}{get.Rul.history}}. The
averages and the drawdown use the whole history, so pass all of it and
zoom with the range buttons.

\item[\code{mode}] \code{"light"} (the default) or \code{"dark"}, to suit a light or dark
page.

\item[\code{averages}] Days of the moving averages to draw (see
\code{\LinkA{moving\_average()}{moving.Rul.average}}), or \code{NULL} for none.

\item[\code{show\_drawdown}] \code{TRUE} to add the drawdown chart (see \code{\LinkA{drawdown()}{drawdown}}).
\end{ldescription}
\end{Arguments}
%
\begin{Value}
A plotly object.
\end{Value}
%
\begin{Examples}
\begin{ExampleCode}
## Not run: 
btc <- get_history("Bitcoin (BTC)", as.Date("2013-04-29"))
plot_timeseries(btc)

# shorter averages, and no drawdown chart
plot_timeseries(btc, averages = c(20, 100), show_drawdown = FALSE)

## End(Not run)
\end{ExampleCode}
\end{Examples}
\HeaderA{plot\_treemap}{Treemap of the top 100 coins}{plot.Rul.treemap}
%
\begin{Description}
Each box is sized by market cap and colored from red to green by the 24h
price change. The color limits are the 5th and 95th percentiles of the
changes, so a few outliers do not wash out the rest. Coins without a 24h
change are labelled n/a and colored white.
\end{Description}
%
\begin{Usage}
\begin{verbatim}
plot_treemap(listings, mode = "light")
\end{verbatim}
\end{Usage}
%
\begin{Arguments}
\begin{ldescription}
\item[\code{listings}] The \code{data} of \code{\LinkA{get\_listings()}{get.Rul.listings}}.

\item[\code{mode}] \code{"light"} (the default) or \code{"dark"}, to suit a light or dark
page.
\end{ldescription}
\end{Arguments}
%
\begin{Value}
A plotly object.
\end{Value}
%
\begin{Examples}
\begin{ExampleCode}
## Not run: 
plot_treemap(get_listings()$data)

## End(Not run)
\end{ExampleCode}
\end{Examples}
\printindex{}
\end{document}
